\scanpage{71}
\margin{\deleted{(\unsure{2 mots})}}
\mspage{11}
\[
\text{(34)}\qquad
\mathcal X(\mathbb P^1_{\mathbb C},\nu_*)=
\deleted{\text{catégorie}}\inserted{\text{topos}}\ \text{des }\deleted{\text{espaces}}\inserted{\text{surfaces}}\ \text{topologiques (\unsure{pas}}
\]
néc. séparées) ramifiées sur $\mathbb P^1_{\mathbb C}$, avec indices de ramification, \deleted{en} les points au dessus d'un $R_i\in\operatorname{supp}\nu_*$, divisant $\nu(R_i)=\nu_i$.

\margin{\uline{NB} Ce topos est une multiplicité topologique, loc. surface topologique, elle est canoniquement orientée.}

i.e. on a un isom. canonique
\[
\text{(35)}\qquad
\Pi_{03}/\langle(\rho_i^{\nu_i})\rangle
\simeq\Pi_1(\mathbb P^1_{(\mathbb C)},\nu_*;P_+)
\overset{\text{déf}}{=}\Pi_1(\mathcal X(\mathbb P^1(\mathbb C),\nu_*),P_+),
\]
et de même
\[
\begin{aligned}
\text{(36)}\qquad
\Pi_{0,3}^{\tau}/\unsure{\langle}(\rho_i^{\nu_i})\rangle
&\simeq\Pi_1(\mathbb P^1(\mathbb C),\{1,\tau\},\nu_*;P_+)\overset{\text{déf}}{=}\\
&\Pi_1(\mathcal X(\mathbb P^1(\mathbb C),\nu_*\deleted{\unsure{\}}},\{1,\tau\}),P_+)
\end{aligned}
\]
où
\[
\text{(37)}\qquad
\mathcal X(\mathbb P^1(\mathbb C),\nu_*,\{1,\tau\})
=(\mathcal X(\mathbb P^1(\mathbb C),\nu_*),\{1,\tau\})
\]
topos mixte défini par le topos (34) avec l'opération naturelle de $\{1,\tau\}$ dessus.

\margin{\uline{NB} c'est encore une multiplicité topologique, localement une surface, mais cette multiplicité n'est pas orientable...}

Passant aux complétés profinis, on trouve
\[
\begin{aligned}
\text{(38)}\qquad
\widehat\Pi_{0,3}/\langle(\rho_i^{\nu_i})\rangle
&\simeq\Pi_1(\underbrace{\mathbb P^1_{\mathbb C}}_{\uline{\text{schéma}}},\nu_*;P_+)\\
&\simeq\Pi_1(\mathbb P^1_{\overline{\mathbb Q}},\nu_*;P_+)
\end{aligned}
\]
\[
\begin{aligned}
\text{(39)}\qquad
\widehat\Pi_{0,3}^{\tau}/\langle(\rho_i^{\nu_i})\rangle
&\simeq\Pi_1(\underbrace{\mathbb P^1_{\mathbb R}}_{\text{schéma}},\nu_*;P_+)\\
&\simeq\Pi_1(\mathbb P^1_{\overline{\mathbb Q}\cap\mathbb R},\nu_*;P_+)
\end{aligned}
\]
(et \deleted{cette} l'interprétation \inserted{(38)} montre qu'il y a une
\margin{70}
